% GENERATED FILE. Edit canonical lessons, then run npm run generate:books.
% canonical-source-hash: fnv1a64:b5e6d8b1e1bb4342
% canonical-lessons: AR-C39-laysa, AR-C39-ma, AR-C39-la-wala, AR-C39-hal, AR-R39-denying-and-asking

\chapter{Denying, and Asking}
\label{ch:ar-denying-and-asking}
\hlchaptermodality{39}

\begin{chapteropening}
\textbf{By the end of this chapter:} \emph{I can deny a nominal sentence with \ar{ليس}, a past-tense verb with \ar{ما} and a present-tense verb with \ar{لا}, refuse two things at once with \ar{لا} … \ar{ولا}, and turn any of them into a yes-or-no question with \ar{هل}.}

Three negators sorted by the kind of sentence each one owns, then every one of them asked back with \ar{هل} -- so \ar{نعم} and \ar{لا}, taught in chapter 5 as answers, finally have a question to answer.
\end{chapteropening}

\section[laysa]{\ar{ليس} (laysa) — is not}

\label{lesson:AR-C39-laysa}



\begin{quote}
\textbf{\ar{اسمي} \ar{حسن}} — \emph{my name is Hasan} — is a complete Arabic sentence with no verb in it, and it is the sentence shape this book teaches best. Try to deny it with \textbf{\ar{لا}} and there is no verb for \textbf{\ar{لا}} to stand in front of.

\begin{itemize}
\raggedright
  \item \emph{Recall:} the joiner two lessons back that has \textbf{\ar{لا}} hiding inside it — \emph{lākin}
\end{itemize}
\end{quote}

\subsection*{You'll want to know: \ar{ليس}}
\begin{quote}
\textbf{\ar{اسمي} \ar{ليس} \ar{حسن}} — \emph{ismī laysa Ḥasan} — \textbf{my name is not Hasan.}
\end{quote}

\begin{quote}
\textbf{\ar{الحليب} \ar{ليس} \ar{ماءً}} — \emph{al-ḥalīb laysa māʾan} — \textbf{milk is not water.}
\end{quote}

Arabic leaves the present-tense \emph{is} out. To deny a sentence that never had one, the language supplies a word that exists only for that job — and \textbf{\ar{ليس}} is it.

Notice the \textbf{\ar{ـً}} on \textbf{\ar{ماءً}}: \textbf{\ar{ليس}} puts what follows it in the accusative, which is the one grammatical demand it makes and the reason its predicate wears a tanwīn where the affirmative would not.

\begin{cousinweb}
Arab grammarians read \textbf{\ar{ليس}} as \textbf{\ar{لا}} — the negator of chapter five — welded to \textbf{\ar{أيس}}, an old word meaning \emph{there is}. So the word began as \emph{there-is-not}, and it still behaves like a verb: it takes person endings, even though nothing it denies is a verb. It is a fossil doing a job nothing else in the language does.
\end{cousinweb}

\subsection*{Your turn}
Say these aloud:

\begin{itemize}
\raggedright
  \item the affirmative, then the denial — \emph{ismī Ḥasan} … \emph{ismī laysa Ḥasan}
  \item \emph{al-ḥalīb laysa māʾan}
  \item \emph{ash-shāy laysa ḥalīban}
\end{itemize}

\begin{tcolorbox}[breakable,colback=teal!4,colframe=teal!35!black,title={Before you move on}]
Why can a nominal sentence not be denied with \textbf{\ar{لا}}? (\textbf{There is no verb for \ar{لا} to stand in front of} — Arabic leaves \emph{is} out.) What are the two pieces grammarians find inside \textbf{\ar{ليس}}? (\textbf{\ar{لا}} and \textbf{\ar{أيس}}, \emph{there is}.) And what does \textbf{\ar{ليس}} do to the word after it? (\textbf{Puts it in the accusative}.)

Source: \href{https://www.unicode.org/charts/PDF/U0600.pdf}{Unicode Arabic chart}.
\end{tcolorbox}
\section[mā]{\ar{ما} (mā) — did not}

\label{lesson:AR-C39-ma}



\begin{quote}
\textbf{\ar{لا}} denies a present-tense verb and \textbf{\ar{ليس}} denies a sentence with no verb at all. Eleven verbs in this book are in the \textbf{past}, and neither of those two touches them.

\begin{itemize}
\raggedright
  \item \emph{Recall:} the item before this one, and what it demands of the word after it — \emph{laysa}
\end{itemize}
\end{quote}

\subsection*{You'll want to know: \ar{ما}}
\begin{quote}
\textbf{\ar{ما} \ar{كتب}} — \emph{mā kataba} — \textbf{he did not write.}
\end{quote}

\begin{quote}
\textbf{\ar{ما} \ar{أكل}} — \emph{mā akala} — \textbf{he did not eat.}
\end{quote}

Arabic uses a \textbf{different word} to negate the past from the one it uses for the present. \textbf{\ar{لا}} in front of a past-tense verb is not a weaker sentence; it is the wrong word.

\begin{grammarlens}[title={Grammar Lens: three negators, three kinds of sentence}]
\noindent
\begin{tabularx}{\linewidth}{@{}>{\raggedright\arraybackslash}X>{\raggedright\arraybackslash}X>{\raggedright\arraybackslash}X@{}}
\toprule
\textbf{negator} & \textbf{denies} & \textbf{example} \\
\midrule
\textbf{\ar{لا}} & a present-tense verb & \emph{lā aktub} \\
\textbf{\ar{ما}} & a past-tense verb & \emph{mā kataba} \\
\textbf{\ar{ليس}} & a sentence with no verb & \emph{ismī laysa Ḥasan} \\
\bottomrule
\end{tabularx}

That is the whole of A1 Arabic negation in three rows, and this chapter is where the second and third of them arrive.
\end{grammarlens}

\begin{cousinweb}
You have written these two letters before. \textbf{\ar{ما} \ar{اسمك}} — \emph{what is your name?} — opens chapter two, and that \textbf{\ar{ما}} means \emph{what}. It is the same word. Semitic languages carry one very general particle that Arabic puts to work as \emph{what}, as \emph{not}, and as \emph{as long as}, and only its position in the sentence tells you which one you are reading.
\end{cousinweb}

\subsection*{Your turn}
Say these aloud:

\begin{itemize}
\raggedright
  \item \emph{kataba} … \emph{mā kataba}
  \item the present beside the past — \emph{lā aktub} … \emph{mā kataba}
  \item the two jobs of the same two letters — \emph{mā ismuka?} … \emph{mā akala}
\end{itemize}

\begin{tcolorbox}[breakable,colback=teal!4,colframe=teal!35!black,title={Before you move on}]
Which negator goes with a past-tense verb, which with a present, and which with no verb? (\textbf{\ar{ما}}, \textbf{\ar{لا}}, \textbf{\ar{ليس}}.) Which question word is spelled the same as this one? (\textbf{\ar{ما}}, \emph{what}.) And what tells the two apart? (\textbf{Where they stand in the sentence}.)

Source: \href{https://www.unicode.org/charts/PDF/U0600.pdf}{Unicode Arabic chart}.
\end{tcolorbox}
\section[lā … wa-lā …]{\ar{لا} … \ar{ولا} … — neither, nor}

\label{lesson:AR-C39-la-wala}



\begin{quote}
The last chapter taught the joiner that adds. This lesson uses it to refuse — which is the one place a coordinator and a negator arrive as a single written word.

\begin{itemize}
\raggedright
  \item \emph{Recall:} the negator of the past, from the lesson before — \emph{mā}
\end{itemize}
\end{quote}

\subsection*{You'll want to know: \ar{لا} … \ar{ولا} …}
\begin{quote}
\textbf{\ar{لا} \ar{شاي} \ar{ولا} \ar{حليب}} — \emph{lā shāy wa-lā ḥalīb} — \textbf{neither tea nor milk.}
\end{quote}

One \textbf{\ar{لا}} in front of the first thing, and \textbf{\ar{ولا}} — the joiner \textbf{\ar{و}} written onto a second \textbf{\ar{لا}} — in front of the next. English needs two different words, \emph{neither} and \emph{nor}; Arabic uses the same word twice and glues the \emph{and} onto the second.

\begin{grammarlens}[title={Grammar Lens: the mirror of the joiner}]
\noindent
\begin{tabularx}{\linewidth}{@{}>{\raggedright\arraybackslash}X>{\raggedright\arraybackslash}X@{}}
\toprule
\textbf{said} & \textbf{means} \\
\midrule
\textbf{\ar{الشاي} \ar{والحليب}} & the tea and the milk \\
\textbf{\ar{لا} \ar{شاي} \ar{ولا} \ar{حليب}} & neither tea nor milk \\
\bottomrule
\end{tabularx}

The second line is the first with a \textbf{\ar{لا}} put in front of each thing. Nothing else changes — and the definite article drops, because you are refusing tea in general rather than a particular cup of it.
\end{grammarlens}

\subsection*{Your turn}
Say these aloud:

\begin{itemize}
\raggedright
  \item \emph{lā shāy wa-lā ḥalīb}
  \item the affirmative it mirrors — \emph{ash-shāy wa-l-ḥalīb}
  \item \emph{lā qahwa wa-lā sukkar}
\end{itemize}

\begin{tcolorbox}[breakable,colback=teal!4,colframe=teal!35!black,title={Before you move on}]
How many negators does \emph{neither … nor} need in Arabic? (\textbf{Two, and they are the same word}.) What is welded onto the second one? (\textbf{\ar{و}}, the joiner.) And what happens to the definite article? (\textbf{It drops} — the refusal is general.)

Source: \href{https://www.unicode.org/charts/PDF/U0600.pdf}{Unicode Arabic chart}.
\end{tcolorbox}
\section[hal]{\ar{هل} (hal) — a question is coming}

\label{lesson:AR-C39-hal}



\begin{quote}
\textbf{\ar{نعم}} and \textbf{\ar{لا}} were taught in chapter five as answers. Thirty-four chapters later this book still cannot ask the question they answer.

\begin{itemize}
\raggedright
  \item \emph{Recall:} the double refusal from the lesson before — \emph{lā shāy wa-lā ḥalīb}
\end{itemize}
\end{quote}

\subsection*{You'll want to know: \ar{هل}}
\begin{quote}
\textbf{\ar{هل} \ar{أنت} \ar{بخير}} — \emph{hal anta bi-khayr?} — \textbf{are you well?}
\end{quote}

\begin{quote}
\textbf{\ar{هل} \ar{كتب}} — \emph{hal kataba?} — \textbf{did he write?}
\end{quote}

Put \textbf{\ar{هل}} at the front of any statement and it becomes a yes-or-no question. Nothing else in the sentence moves — no word order changes, no ending changes.

Arabic \textbf{needs} the word. English asks by moving the verb in front of the subject; Bengali and Urdu ask by letting the voice rise. Written Arabic does neither: the particle carries the whole question, which is why a page of Arabic tells you a question is coming before you have read any of it.

\begin{grammarlens}[title={Grammar Lens: what the particle replaces}]
\noindent
\begin{tabularx}{\linewidth}{@{}>{\raggedright\arraybackslash}X>{\raggedright\arraybackslash}X@{}}
\toprule
\textbf{language} & \textbf{how a yes-or-no question is marked} \\
\midrule
English & the verb moves in front of the subject \\
Arabic & \textbf{\ar{هل}} at the front, and nothing moves \\
\bottomrule
\end{tabularx}

A negative asks the same way, and that is the ordinary way to check something you suspect: \emph{hal laysa hunāka shāy?} Answer either with \textbf{\ar{نعم}} or with \textbf{\ar{لا}}.
\end{grammarlens}

\subsection*{Your turn}
Say these aloud:

\begin{itemize}
\raggedright
  \item \emph{hal anta bi-khayr?} — and answer it \emph{naʿam} or \emph{lā}
  \item \emph{hal kataba?}
  \item the statement, then the question — \emph{anta bi-khayr} … \emph{hal anta bi-khayr?}
\end{itemize}

\begin{tcolorbox}[breakable,colback=teal!4,colframe=teal!35!black,title={Before you move on}]
What moves when an Arabic statement becomes a yes-or-no question? (\textbf{Nothing} — \textbf{\ar{هل}} goes at the front.) Why can it not be done with the voice alone, as in Bengali or Urdu? (\textbf{Written Arabic marks it with a particle}.) And which two words answer it? (\textbf{\ar{نعم}} and \textbf{\ar{لا}}, from chapter five.)

Source: \href{https://www.unicode.org/charts/PDF/U0600.pdf}{Unicode Arabic chart}.
\end{tcolorbox}
\section[laysa · mā · hal]{Denying, and asking}

\label{lesson:AR-R39-denying-and-asking}



\begin{quote}
Arabic has three negators and each one owns a different kind of sentence. Name all three, and say what each one denies.
\end{quote}

\subsection*{Your turn}
\begin{itemize}
\raggedright
  \item \emph{Say it:} the verbless sentence denied — \emph{ismī laysa Ḥasan}
  \item \emph{Say it:} the past denied — \emph{mā kataba}
  \item \emph{Say it:} the present denied — \emph{lā aktub}, and \emph{lā afham}
  \item \emph{Say it:} both refused at once — \emph{lā shāy wa-lā ḥalīb}
  \item \emph{Say it:} each of those asked back — put \emph{hal} in front and change nothing else
  \item \emph{Recall:} six lessons back, the word that accepts what was said — \emph{ḥasanan}
\end{itemize}

\begin{tcolorbox}[breakable,colback=teal!4,colframe=teal!35!black,title={Before you move on}]
Which negator would you use for \textbf{\ar{اسمي} \ar{حسن}}, which for \textbf{\ar{كتب}}, and which for \textbf{\ar{أكتب}}? (\textbf{\ar{ليس}}, then \textbf{\ar{ما}}, then \textbf{\ar{لا}}.) What does \textbf{\ar{هل}} ask of the sentence it fronts? (\textbf{Nothing at all} — it only marks it.)
\end{tcolorbox}
